\documentclass[11pt,a4paper]{article}

% ── Packages ──────────────────────────────────────────────────────────────────
\usepackage[utf8]{inputenc}
\usepackage[T1]{fontenc}
\usepackage[margin=1in]{geometry}
\usepackage{amsmath,amssymb}
\usepackage{booktabs}
\usepackage{graphicx}
\usepackage{hyperref}
\usepackage{caption}
\usepackage{subcaption}
\usepackage{multirow}
\usepackage{xcolor}
\usepackage{enumitem}
\usepackage{float}
\usepackage[numbers]{natbib}

\hypersetup{
    colorlinks=true,
    linkcolor=blue!70!black,
    citecolor=green!50!black,
    urlcolor=blue!60!black,
}

% ── Title ─────────────────────────────────────────────────────────────────────
\title{\textbf{Multi-Model Fusion for Robust 2D Upper-Limb Pose Estimation}}
\author{
    Jerry Chen \\
    \textit{University of Pittsburgh} \\
    \texttt{jzc23@pitt.edu}
}
\date{April 2026}

\begin{document}
\maketitle

% ══════════════════════════════════════════════════════════════════════════════
\begin{abstract}
Single pose-estimation models, while individually accurate, are prone to
sporadic catastrophic failures caused by occlusion, motion blur, or
domain-specific appearance changes.  We present a learned fusion framework
that combines the predictions of three complementary 2D pose estimators---HRNet/DEKR,
OpenPose, and MediaPipe---to produce robust upper-limb keypoint estimates for
exercise and rehabilitation video.  We compare an MLP-based attention fusion, a
Transformer encoder--decoder fusion, an internal-attention Transformer variant,
and two rule-based baselines (uniform and confidence-weighted averaging) on a
custom dataset of 5{,}346 test frames across three left-arm joints (16{,}038
evaluation samples).  The MLP Fusion achieves \textbf{100\% accuracy} at a
normalized distance threshold of $\tau{=}0.05$ ($\approx$18\,px), eliminating
all outlier predictions while maintaining a mean L2 error of 3.22\,px.  We
discuss the \emph{precision--robustness tradeoff}: HRNet alone attains a lower
mean L2 (1.80\,px) yet fails on 0.34\% of samples, whereas the fusion model
hedges its predictions to avoid every failure.  Code and evaluation pipelines
are publicly available.
\end{abstract}

% ══════════════════════════════════════════════════════════════════════════════
\section{Introduction}\label{sec:intro}

Accurate 2D human pose estimation is a foundational task in computer vision
with direct applications in physical therapy, exercise monitoring, and
biomechanical analysis.  Modern deep-learning--based pose estimators such as
HRNet/DEKR~\cite{geng2021dekr, sun2019hrnet}, OpenPose~\cite{cao2017openpose},
and MediaPipe~\cite{lugaresi2019mediapipe} achieve high average accuracy, yet
each model exhibits distinct failure modes: HRNet can mislocalize joints during
fast motion, OpenPose may miss detections entirely under occlusion, and
MediaPipe occasionally produces large outliers on uncommon viewpoints.

In safety-sensitive applications---for instance, tracking a patient's arm
during rehabilitation exercises---a single catastrophic misprediction can
invalidate an entire session.  Rather than investing in a single
ever-larger model, we explore \emph{multi-model fusion}: can a lightweight
learned combiner, trained on the precomputed outputs of several off-the-shelf
estimators, eliminate outlier failures while preserving overall precision?

\paragraph{Contributions.}
\begin{enumerate}[nosep]
    \item A per-joint attention-based MLP fusion architecture that achieves
          100\% threshold-based accuracy on a rehabilitation-video test set,
          surpassing every individual estimator.
    \item A systematic comparison of learned (MLP, Transformer,
          Internal-Attention Transformer) and rule-based (uniform,
          confidence-weighted) fusion strategies.
    \item An analysis of the \emph{precision vs.\ robustness tradeoff},
          explaining why fusion models may have higher mean error yet fewer
          catastrophic failures than the best single estimator.
    \item A study of fine-tuning strategies for pose backbones (full
          fine-tuning vs.\ frozen backbone), including diagnosis and resolution
          of catastrophic forgetting.
\end{enumerate}

% ══════════════════════════════════════════════════════════════════════════════
\section{Related Work}\label{sec:related}

\paragraph{Pose estimation backbones.}
Bottom-up approaches such as DEKR~\cite{geng2021dekr} built on
HRNet~\cite{sun2019hrnet} regress disentangled per-keypoint heatmaps and
offsets, achieving strong multi-person results on COCO~\cite{lin2014coco}.
OpenPose~\cite{cao2017openpose} uses Part Affinity Fields for real-time
multi-person detection.  MediaPipe~\cite{lugaresi2019mediapipe} targets
mobile and edge deployment with a lightweight pipeline.
ViTPose~\cite{xu2022vitpose} applies Vision Transformers to pose estimation,
offering state-of-the-art results on multiple benchmarks.

\paragraph{Model fusion and ensembles.}
Ensemble methods are well-established in classification~\cite{dietterich2000ensemble}
but less explored for structured regression outputs such as keypoint
coordinates.  Prior work on pose ensemble typically averages heatmaps before
decoding~\cite{chen2018cascaded}, whereas our approach operates directly on
decoded coordinate--confidence tuples, enabling heterogeneous model
combinations without requiring access to internal feature maps.

\paragraph{Attention-based fusion.}
Multi-head self-attention~\cite{vaswani2017attention} provides a natural
mechanism for weighting variable-quality inputs per token.  Our MLP fusion
treats each model's prediction for a given joint as a token and applies a
single layer of self-attention to produce a fused output---a lightweight
alternative to full Transformer encoder--decoder architectures.

% ══════════════════════════════════════════════════════════════════════════════
\section{Dataset}\label{sec:dataset}

\paragraph{Video corpus.}
We use a custom multi-camera dataset of exercise and rehabilitation sessions
captured at two camera angles (\texttt{cam0}, \texttt{cam2}).  Video clips are
organized by patient and exercise type under a hierarchical directory
(\texttt{clips\_mp4/<patient>/<exercise>/}).  Each clip is stored as an MP4
file with accompanying ground-truth annotations.

% >>> AUTHOR NOTE: replace or expand the following sentence with the
%     institutional source, IRB approval, or DOI if applicable. <<<
\textit{(Dataset source: collected from rehabilitation/exercise sessions at the
University of Pittsburgh.  Details on collection protocol and IRB approval
should be added here.)}

\paragraph{Annotations.}
Ground-truth 2D keypoints are provided in per-video \texttt{p2d} text files.
Each line contains space-separated floating-point coordinates for the three
left-arm joints of interest: \textbf{left shoulder}, \textbf{left elbow}, and
\textbf{left wrist} (6 values: $x_1, y_1, x_2, y_2, x_3, y_3$).

\paragraph{Data splits.}
Splits are performed at the patient level to prevent data leakage.  The test
partition comprises \textbf{5{,}346 frames} across multiple videos.  Since each
frame contributes three joint evaluations, the total number of evaluation
samples is $5{,}346 \times 3 = \mathbf{16{,}038}$.

\paragraph{Precomputed prediction caches.}
Rather than running each backbone at training time, we precompute per-frame
2D predictions and confidence scores for every model (HRNet/DEKR, OpenPose,
MediaPipe) and store them in cache files under \texttt{p2d\_cache/}.  The
fusion models are trained and evaluated entirely on these cached predictions,
making the pipeline fast and reproducible.

% ══════════════════════════════════════════════════════════════════════════════
\section{Methods}\label{sec:methods}

All fusion models receive the same input per frame: for each of $K{=}3$ joints,
the $(x,y)$ coordinates and scalar confidence from each of $M{=}3$ backbone
models, yielding a tensor of shape $(B, K, 3)$ per model (2 coordinates + 1
confidence).

\subsection{Baseline Detectors}\label{sec:baselines}

\begin{description}[nosep, leftmargin=1em]
    \item[HRNet/DEKR] Bottom-up keypoint regression using High-Resolution
        Network (HRNet-W48) with Disentangled Keypoint Regression, pretrained
        on COCO~\cite{geng2021dekr, sun2019hrnet}.
    \item[OpenPose] Real-time multi-person pose estimation via Part Affinity
        Fields, using the CMU BODY\_25 model~\cite{cao2017openpose}.
    \item[MediaPipe] Google's lightweight pose-estimation pipeline optimized
        for mobile inference~\cite{lugaresi2019mediapipe}.
\end{description}

\subsection{Fusion Architectures}\label{sec:fusion-arch}

\subsubsection{MLP Attention Fusion}

For each joint $k$, the three model predictions (each a 3-vector of $[x, y, \text{conf}]$) are
independently projected to a $d$-dimensional embedding via per-model \texttt{SmallMLP}
modules, then summed with learned joint and model embeddings.  A single
multi-head self-attention layer ($d{=}64$, 4 heads) fuses the three model
tokens into a weighted combination, followed by a feed-forward network and
layer normalization.  Two linear heads decode the fused embedding into
predicted coordinates $(x, y)$ and a confidence logit.  The attention weights
provide interpretable \emph{importance scores} per model per joint.

\subsubsection{Transformer Encoder--Decoder Fusion}

Each model's per-joint features are encoded by a dedicated
\texttt{TransformerEncoderLayer} and concatenated into a memory of length
$3K$.  A learned query set of $K$ tokens is decoded via cross-attention
(\texttt{FusionTransformerDecoderLayer}) to produce the fused
output~\cite{vaswani2017attention}.  The cross-attention weights serve as
importance indicators.

\subsubsection{Internal-Attention Transformer Fusion}

A self-attention variant where all $3K$ model--joint tokens are processed by
stacked Transformer encoder layers with $3K$-length self-attention.  A
pooling-and-reshape step maps the output back to $K$ joint predictions.

\subsection{Rule-Based Baselines}\label{sec:rule-based}

\begin{description}[nosep, leftmargin=1em]
    \item[Uniform Fusion] Each model receives equal weight
        $w_m = \tfrac{1}{3}$; the fused coordinate is the arithmetic mean.
    \item[Confidence-Weighted Fusion] Weights are proportional to each
        model's reported confidence per joint:
        $w_m^{(k)} = c_m^{(k)} / \sum_{m'} c_{m'}^{(k)}$.
\end{description}

\subsection{Fine-Tuned Single Models}\label{sec:finetuned}

To contextualize fusion performance, we also fine-tune individual backbones
on the target domain:

\begin{itemize}[nosep]
    \item \textbf{DEKR (trained)}: Full fine-tuning from COCO weights with a
        conservative learning rate ($10^{-5}$) and ColorJitter augmentation.
    \item \textbf{DEKR (frozen backbone)}: Only the regression head is
        trained; backbone weights are frozen.
    \item \textbf{OpenPose (trained / frozen)}: Analogous strategies applied
        to the OpenPose model.
\end{itemize}

Key training insights from iterative experimentation (see
Section~\ref{sec:discussion}): aggressive learning rates caused catastrophic
forgetting in DEKR (accuracy collapsed to ${\sim}10\%$); a fresh restart with
reduced LR and augmentation recovered performance to 99.95\%.

\subsection{Training Details}

All fusion models are trained with the AdamW optimizer on the cached
three-model predictions.  The loss combines a smooth-L1 coordinate regression
term with a binary cross-entropy confidence loss.  Data augmentation for fusion
models includes random keypoint jitter ($\pm 2$\,px) and 10\% random joint
occlusion to improve robustness.

% ══════════════════════════════════════════════════════════════════════════════
\section{Experiments and Results}\label{sec:results}

\subsection{Evaluation Protocol}

We evaluate on the held-out test set of 16{,}038 joint predictions (5{,}346
frames $\times$ 3 joints).  Metrics:

\begin{itemize}[nosep]
    \item \textbf{Accuracy} ($\tau{=}0.05$): fraction of predictions within a
        normalized L2 distance of 0.05 from ground truth ($\approx$18\,px in
        $256{\times}256$ image space).
    \item \textbf{Mean L2 (px)}: average Euclidean error in pixels.
    \item \textbf{Mean L2 (norm)}: L2 normalized by image diagonal.
    \item \textbf{MAE}: mean absolute error of coordinates.
    \item \textbf{Calibration gap}: difference between mean confidence and
        actual accuracy (positive = overconfident).
\end{itemize}

\subsection{Fusion Model Comparison}

Table~\ref{tab:fusion} compares the five fusion strategies.  MLP Fusion achieves
perfect accuracy while Confidence-Weighted Fusion attains the lowest mean L2.

\begin{table}[H]
\centering
\caption{Fusion model comparison on the test set (16{,}038 samples).
         Best value in each column is \textbf{bolded}.}
\label{tab:fusion}
\begin{tabular}{@{}lcccc@{}}
\toprule
\textbf{Model} & \textbf{Accuracy (\%)} & \textbf{L2 (px)} & \textbf{L2 (norm)} & \textbf{MAE} \\
\midrule
MLP Fusion                    & \textbf{100.00} & 3.215 & 0.0089 & 0.0076 \\
Confidence-Weighted           & 99.92  & \textbf{2.074} & \textbf{0.0057} & \textbf{0.0047} \\
Transformer Fusion            & 99.69  & 4.492 & 0.0124 & 0.0114 \\
Internal Transformer Fusion   & 99.58  & 3.441 & 0.0095 & 0.0085 \\
Uniform Fusion ($\tfrac{1}{3}$) & 98.39  & 3.052 & 0.0084 & 0.0073 \\
\bottomrule
\end{tabular}
\end{table}

\subsection{Individual Baseline Comparison}

Table~\ref{tab:baselines} reports the performance of each backbone model
individually, including fine-tuned variants.

\begin{table}[H]
\centering
\caption{Individual baseline model performance.  ``Pre-trained'' models use
         off-the-shelf COCO/ImageNet weights; ``Trained'' models are fine-tuned
         on the target domain.}
\label{tab:baselines}
\begin{tabular}{@{}lccc@{}}
\toprule
\textbf{Model} & \textbf{Accuracy (\%)} & \textbf{L2 (px)} & \textbf{Samples} \\
\midrule
\multicolumn{4}{@{}l}{\textit{Pre-trained (cached predictions)}} \\
\quad HRNet/DEKR              & 99.66 & 1.795  & 16{,}038 \\
\quad MediaPipe               & 98.74 & 2.532  & 15{,}972 \\
\quad OpenPose                & 98.18 & 2.612  & 14{,}916 \\
\midrule
\multicolumn{4}{@{}l}{\textit{Fine-tuned on target domain}} \\
\quad DEKR (trained)          & 99.95 & 6.273  & 16{,}038 \\
\quad OpenPose (frozen backbone) & 91.71 & 9.778  & 16{,}038 \\
\quad OpenPose (trained)      & 88.50 & 11.378 & 16{,}038 \\
\quad DEKR (frozen backbone)  & 70.72 & 29.733 & 16{,}038 \\
\bottomrule
\end{tabular}
\end{table}

\subsection{Patient-Level Model Evaluation}

To complement the fusion-specific test set evaluation, we conducted a broader
frame-level analysis of five backbone models across the full dataset
(52 patient--camera cases, 47{,}082 annotated frames).  For each model, the
normalized mean error per frame was computed, then aggregated as a
frame-weighted mean with 95\% confidence intervals.  Table~\ref{tab:patient}
summarizes these results.

\begin{table}[H]
\centering
\caption{Frame-level mean error (normalized) across 52 patient--camera cases
         (47{,}082 frames).  The weighted mean accounts for varying video
         lengths; 95\% CI is computed from the pooled standard deviation.}
\label{tab:patient}
\begin{tabular}{@{}lccccc@{}}
\toprule
\textbf{Model} & \textbf{Weighted Mean} & \textbf{95\% CI} & \textbf{Pooled Std} & \textbf{Unweighted Mean} & \textbf{Cases} \\
\midrule
ViTPose    & 0.0302 & [0.0301, 0.0303] & 0.0128 & 0.0308 & 52 \\
DEKR       & 0.0337 & [0.0335, 0.0338] & 0.0163 & 0.0341 & 52 \\
OpenPose   & 0.0354 & [0.0352, 0.0356] & 0.0234 & 0.0369 & 52 \\
MediaPipe  & 0.0357 & [0.0356, 0.0359] & 0.0148 & 0.0366 & 52 \\
PoseBERT   & 0.0669 & [0.0666, 0.0672] & 0.0368 & 0.0661 & 52 \\
\bottomrule
\end{tabular}
\end{table}

ViTPose achieves the lowest frame-level error (0.0302), followed closely by
DEKR (0.0337), OpenPose (0.0354), and MediaPipe (0.0357).  PoseBERT, designed
primarily for 3D mesh recovery rather than precise 2D joint localization,
trails at 0.0669.  The tight confidence intervals (owing to the large frame
count) confirm that the ranking is statistically robust.  Notably, the relative
ordering differs from the fusion test set results (Table~\ref{tab:baselines}),
where DEKR/HRNet leads in L2; this is because the patient-level evaluation
spans the full dataset rather than the held-out test split and uses a different
error normalization.

\subsection{Confidence Calibration}

Table~\ref{tab:calibration} shows the confidence calibration of the learned fusion
models.  MLP Fusion exhibits near-perfect calibration, while the Transformer
variants are slightly overconfident.

\begin{table}[H]
\centering
\caption{Confidence calibration of learned fusion models.  Calibration gap =
         mean confidence $-$ accuracy; negative values indicate
         underconfidence.}
\label{tab:calibration}
\begin{tabular}{@{}lcccc@{}}
\toprule
\textbf{Model} & \textbf{Mean Conf.} & \textbf{Conf.\ (correct)} & \textbf{Conf.\ (incorrect)} & \textbf{Cal.\ Gap} \\
\midrule
MLP Fusion                  & 0.9946 & 0.9946 & ---    & $-0.005$ \\
Transformer Fusion          & 0.9992 & 0.9993 & 0.9944 & $+0.002$ \\
Internal Transformer        & 0.9997 & 0.9997 & 0.9998 & $+0.004$ \\
Uniform Fusion              & 0.9011 & 0.9056 & 0.6250 & $-0.078$ \\
Confidence-Weighted         & 0.9433 & 0.9433 & 0.9407 & $-0.056$ \\
\bottomrule
\end{tabular}
\end{table}

\subsection{Importance Weight Analysis}

The attention mechanism in learned fusion models produces per-frame, per-joint
\emph{importance weights} indicating how much each backbone contributes to the
fused output.  Analysis of these weights across the 5{,}346 test frames reveals
that HRNet/DEKR consistently receives the highest importance in the MLP model,
followed by MediaPipe and OpenPose.  Notably, the weights shift at video
boundaries and during challenging frames (occlusion, fast motion), demonstrating
that the model dynamically adapts its weighting strategy.

% ══════════════════════════════════════════════════════════════════════════════
\section{Discussion}\label{sec:discussion}

\subsection{Precision vs.\ Robustness Tradeoff}

A central finding is the tension between \emph{mean error} and
\emph{worst-case reliability}.  HRNet/DEKR achieves the lowest mean L2
(1.795\,px) among all models but fails the accuracy threshold on 0.34\% of
samples.  MLP Fusion, by contrast, achieves 100\% accuracy---every prediction
falls within $\tau{=}0.05$---at the cost of a higher mean L2 (3.215\,px).

This behavior arises because the fusion model learns to ``hedge'' its
predictions.  When one backbone produces an outlier, the attention mechanism
down-weights it in favor of the remaining models, pulling the prediction back
within the safe threshold.  The result is a small increase in average error
(the fused prediction is pulled away from the best single-model estimate on
easy frames) but a complete elimination of catastrophic failures.

For safety-critical applications such as patient monitoring, this tradeoff is
clearly favorable: practitioners typically set hard error thresholds, and
\emph{zero} threshold violations is more valuable than a marginally lower
average error.

\subsection{MLP vs.\ Transformer Fusion}

Somewhat counter-intuitively, the simplest learned architecture---MLP
Fusion with a single attention layer---outperforms both Transformer variants
on accuracy (100\% vs.\ 99.69\% and 99.58\%).  The Transformers' larger
parameter count increases overfitting risk, despite stronger regularization
(dropout~0.2, keypoint jitter, random occlusion).  The MLP model's inductive
bias of treating each joint independently via per-joint attention appears
well-suited to this low-dimensional fusion task where cross-joint reasoning
offers limited benefit.

\subsection{Confidence-Weighted Fusion: A Strong Heuristic}

The confidence-weighted rule-based baseline achieves the lowest mean L2
(2.074\,px) and near-perfect accuracy (99.92\%).  This suggests that the
backbone models' confidence scores are informative per-sample quality
indicators.  However, the remaining 0.08\% failure rate (13 samples) shows
that confidence alone is insufficient to detect all failure modes---the
learned MLP model captures subtler correlations that confidence cannot.

\subsection{Fine-Tuning Challenges}

Fine-tuning pretrained backbones on the target domain revealed significant
pitfalls:

\begin{itemize}[nosep]
    \item \textbf{Catastrophic forgetting}: DEKR with an aggressive learning
        rate ($5{\times}10^{-5}$) collapsed to ${\sim}10\%$ accuracy.  Reducing
        LR to $10^{-5}$ with data augmentation and a fresh checkpoint
        recovered performance to 99.95\%.
    \item \textbf{Frozen backbone paradox}: OpenPose's frozen-backbone
        variant (91.71\%) \emph{outperformed} its fully-trained counterpart
        (88.50\%), indicating that full fine-tuning suffered from forgetting
        useful COCO features.
    \item \textbf{Strategy mismatch}: Frozen-backbone models require
        \emph{higher} learning rates (the backbone cannot be damaged) and more
        training epochs, the opposite of full fine-tuning best practices.
\end{itemize}

These findings underscore that \textbf{fusion of fixed pretrained models}
can be more practical than investing engineering effort into domain-specific
fine-tuning, which risks degrading already-strong baselines.

\subsection{Limitations}

\begin{itemize}[nosep]
    \item \textbf{Single dataset}: All results are on one exercise/rehabilitation
        corpus; cross-dataset generalization is not evaluated.
    \item \textbf{Three joints only}: The current evaluation targets left-arm
        joints; scaling to full-body pose may change the fusion dynamics.
    \item \textbf{Threshold sensitivity}: The 100\% accuracy claim depends on
        $\tau{=}0.05$; at tighter thresholds, some MLP failures would appear.
    \item \textbf{No temporal modeling}: Each frame is fused independently;
        incorporating temporal consistency could further improve robustness.
    \item \textbf{ViTPose integration}: A fourth backbone (ViTPose~\cite{xu2022vitpose})
        was partially integrated but not included in the final evaluation due
        to coordinate normalization mismatches in the cached predictions.
\end{itemize}

% ══════════════════════════════════════════════════════════════════════════════
\section{Conclusion}\label{sec:conclusion}

We demonstrated that a lightweight learned fusion of three complementary pose
estimators can achieve \textbf{100\% threshold-based accuracy} on a
rehabilitation video test set, entirely eliminating the catastrophic outlier
predictions that affect each individual model.  The key insight is the
\emph{precision--robustness tradeoff}: fusion models sacrifice a small amount
of average accuracy to guarantee that no prediction exceeds the safety
threshold.

Among fusion strategies, the MLP Attention model proved most effective for
accuracy, while the confidence-weighted heuristic offered the best mean L2.
Transformer-based architectures, despite greater capacity, were more prone to
overfitting in this low-dimensional setting.  Our fine-tuning experiments
further showed that domain adaptation of pretrained backbones is fragile and
often counterproductive, reinforcing the case for fusion-level solutions.
\paragraph{Future work.}
Promising directions include (1) integrating ViTPose as a fourth backbone in
the fusion pipeline, (2) adding temporal consistency via recurrent or
Transformer-based sequence models, (3) evaluating on cross-domain benchmarks
(e.g., clinical vs.\ sports video), and (4) exploring the 4-model fusion
variants already prototyped in the codebase.

% ══════════════════════════════════════════════════════════════════════════════
\section{Resources}\label{sec:resources}

\paragraph{Code.}
The full source code for the fusion models, training scripts, evaluation
pipeline, and visualization tools is publicly available at:

\begin{center}
\url{https://github.com/jzc23/pose-fusion}
\end{center}

\noindent\textit{(Update the URL above with the actual public repository link.)}

\paragraph{Dataset.}
The evaluation dataset consists of exercise and rehabilitation video clips
collected at the University of Pittsburgh.  Videos are organized by patient
and exercise type, with manual 2D keypoint annotations for three left-arm
joints.  The dataset structure, split scripts, and annotation format are
documented in the repository.

% >>> AUTHOR NOTE: If the dataset has a DOI, data-sharing agreement, or
%     specific institutional citation, add it here. <<<

\paragraph{Pretrained models and libraries.}
The following open-source projects are used as backbone pose estimators:
\begin{itemize}[nosep]
    \item DEKR / HRNet: \url{https://github.com/HRNet/DEKR}~\cite{geng2021dekr}
    \item OpenPose (PyTorch): \url{https://github.com/CMU-Perceptual-Computing-Lab/openpose}~\cite{cao2017openpose}
    \item MediaPipe: \url{https://github.com/google/mediapipe}~\cite{lugaresi2019mediapipe}
    \item ViTPose: \url{https://github.com/ViTAE-Transformer/ViTPose}~\cite{xu2022vitpose}
    \item easy\_ViTPose wrapper: \url{https://github.com/JunkyByte/easy_ViTPose}
\end{itemize}

% ══════════════════════════════════════════════════════════════════════════════
\bibliographystyle{plainnat}
\bibliography{references}

\end{document}
